\chapter{vimrc Cookbook}
\label{app:vimrc}

Every vimrc line introduced across this book, collected into one place. This is not a single recommended configuration to copy wholesale; it is the complete set of individual settings this book has justified individually, in the chapter noted, for a reader to select from according to their own needs, following the incremental spirit of Chapter~\ref{ch:installing}'s original minimal vimrc.

\begin{lstlisting}
" Core sanity (Chapter 3)
set nocompatible
syntax on
set number
set incsearch
set ignorecase
set smartcase

" Indentation (Chapter 3)
set expandtab
set shiftwidth=4
set softtabstop=4

" Undo persistence across sessions (Chapter 9)
set undofile
" set undodir=~/.vim/undodir    " optional: centralize undo files

" Command-line completion (Chapter 13)
set wildmenu

" gVim interface, if applicable (Chapter 3)
set guioptions-=m              " remove menu bar
set guioptions-=T              " remove toolbar

" Prose and writing (Chapter 20)
" set spell                    " enable per-file with :set spell instead
" set textwidth=80             " set per-filetype rather than globally

" A custom mapping (Chapter 17)
nnoremap <C-s> <C-a>h

" An autocommand (Chapter 17)
augroup PythonIndent
    autocmd!
    autocmd BufWritePre *.py normal! gg=G
augroup END
\end{lstlisting}

The two options left commented out, \texttt{spell} and \texttt{textwidth}, are shown that way deliberately: both are more useful set per-file-type or invoked on demand than enabled globally for every buffer regardless of content, and a reader adapting this cookbook should treat every line here as a starting point for judgment rather than a fixed prescription, consistent with this book's general preference throughout for understanding a setting's effect over copying it uncritically.
